\section{Construction of the map}
\label{section:construction}

\subsubsection{} In this section we will construct the map
  \[\eta_G : \Poinc_! \cdot \PsId_\Whit^{-1} \to \PsId_! \cdot \coeff_*^\vee [2 \delta_G].\]

\subsection{Reminders on Vinberg's semigroup}

\subsubsection{} We write
    \[T_\ad = T/Z_G \simeq (\Gm)^\Delta\]
for the adjoint torus. We write
    \[T_\ad^+ = (\AA^1)^\Delta,\]
viewed as a semigroup. The points in $T_\ad^+$ with zero-one co\"ordinates are in natural bijection with subsets of $\Delta$ and hence standard parabolics. We write $c_P$ for the point corresponding to the parabolic $P$. The $T$ orbit of $c_P$ identifies with $(T_M)_\ad \simeq T/Z_M$, and these orbits form a stratification of $T_\ad^+$.

\subsubsection{Vinberg's semigroup} Consider the group homomorphism
    \[G \overset{Z_G}{\times} T \to T_\ad.\]
This map is naturally $G \times T \times G$ equivariant, where $G \times G$ acts on the left and right on the source and trivially on the target. Vinberg's semigroup $\Vin_G$ is an algebraic semigroup equipped with a $G \times G \times T$-equivariant semigroup morphism
    \[\Vin_G \to T_\ad^+\]
extending the above group homomorphism. This morphism has a canonical ``unit section''
    \[s : T^+_\ad \to \Vin_G\]
whose orbit under $G \times G$ is an open subscheme denoted $\Vin_G^\circ$.

\subsubsection{Stratification} The fiber of
    \[T_\ad^+ \to \Vin_G^\circ \mon \Vin_G\]
over $c_P$ is canonically identified with
\begin{equation}
\label{eqn:Vin_G_strata}    
    \pt \simeq P^-/U_P^- \overset{M}{\times} U_P \bs P
    \to G/U_P^- \overset{M}{\times} U_P \bs G 
    \mon \overline{G/U_P^- \overset{M}{\times} U_P \bs G}
\end{equation}
with the natural $G \times G$ action. The residual action of $Z_M$ is descended from the right and left actions on $G/U_P^-$ and $U_P \bs G$, respectively.

\subsubsection{Generic maps} Let us recall the notion of generic maps from [J. Barlev]. For any scheme $S$, an open subset $U \mon X \times S$ is called a domain if it is open dense after any base change $S' \to S$. For any target $\YY$, we define the prestack
    \[
    \Maps^\gen(X, \YY) : S \mapsto
    \lim_U \Maps(U, \YY)
    \]
where the limit is taken over all domains $U \mon X \times S$. Given a map $\YY_1 \to \YY_2$, we define $\Maps^\gen(X, \YY_1 \to \YY_2)$ as the fiber product:
\[
\begin{tikzcd}
    \Maps^\gen(X, \YY_1 \to \YY_2) \arrow[r] \arrow[d] &  \Maps(X, \YY_2) \arrow[d] \\
    \Maps^\gen(X, \YY_1) \arrow[r] & \Maps^\gen(X, \YY_1)
\end{tikzcd}
\]

\subsubsection{} We define
    \[
    \Vin\Bun_G = \Maps^\gen \left(X, G \bs \Vin^\circ_G / G \mon G \bs \Vin_G / G \right).
    \]
This a stack over $\Bun_G \times \Bun_G \times T_\ad^+$, and we may take the quotient by $T$ to obtain
    \[\Bunc_G = \VBun_G / T,\]
which is a stack over $\Bun_G \times \Bun_G \times T_\ad^+/T$. The stratification of $T_\ad^+/T$ by the poset of standard parabolics induces a stratification of $\Bunc_G$, whose pieces we shall denote $\Bunc_{G,P}$. Note that the open stratum identifies with
    \[\Bunc_{G,G} \simeq \Bun_G \times \BB Z_G.\]

\subsubsection{Reminders on $\PsId_*$ and $\PsId_!$} There is a sequence of maps
  \[\Bun_G \overset{r}{\longrightarrow} \Bun_G \times \BB Z_G \overset{j}{\hookrightarrow} \Bunc_G \overset{\Deltac}{\longrightarrow} \Bun_G \times \Bun_G\]
with the following features:
\begin{itemize}
    \item We have $r_! \simeq r_* [z_G]$.
    \item The map $\Deltac$ is proper.
\end{itemize}
Therefore, we obtain a map
    \[\Delta_! \cdot k_{\Bun_G} \to \Delta_* \cdot k_{\Bun_G} [z_G] \simeq \Delta_* \cdot \omega_{\Bun_G} [z_G - 2 \delta_G].\]
This gives the map
    \[\PsId_! \to \PsId_* [z_G - 2 \delta_G].\]

\subsection{Definition of $\eta_G$}

\subsubsection{Notation} We will write
    \[r_{G,!} = j_! \cdot r_! \cdot k_{\Bun_G} \in \DMod\left( \Bunc_G \right).\]

\subsubsection{} The object $\WW_\vac$ extends cleanly along $\Bun_N \mon \Bunc_N$. We use the same symbol for this extension, and we make the identification
    \[\Whit\left(\Bunc_N\right) \simeq \Whit(\Bun_N).\]
We shall write the definition of $\eta_G$ using $\Bunc_N$ even though it is slightly more complicated because the $\Bunc_N$ version will be needed in the construction of the Jacquet functor.

\subsubsection{} Consider the Cartesian diagram below.
\[
\begin{tikzcd}[
  column sep={4em,between origins},
  row sep={4em,between origins}
]
    & \Bunc_N \arrow[dl] \arrow[dr] & & & \\    
  \Bun_G \arrow[dr] & & \Bunc_N \underset{\Bun_G}{\times} \Bunc_G \arrow[dl] \arrow[dr, "\aph\Deltac"] & & \\
    & \Bunc_G \arrow[dr] & & \Bunc_N \times \Bun_G \arrow[dl] \arrow[dr, "\pr_2"] & \\
    &  & \Bun_G \times \Bun_G & & \Bun_G
\end{tikzcd}
\]
We write
    \[\pi_2 = \pr_2 \cdot \aph\Deltac.\]
Note that we have the identifications
    \[
    \Poinc_! \cdot \WW_\vac \simeq \pi_{2,!} \left( \WW_\vac \stoximes r_{G,!} \right)
    \qquad
    \PsId_! \cdot \coeff_*^\vee \cdot \WW_{\vac,\co} \simeq \pi_{2,*} \left( \WW_\vac \shoximes r_{G,!} \right)
    \]
by base change. Therefore, we expect that the purity morphism and forgetting supports combine to give a map
    \[
    \pi_{2,!} \left( \WW_\vac \shoximes r_{G,!} \right)
    \to \pi_{2,*} \left( \WW_\vac \stoximes r_{G,!} \right) [2 \delta_G].
    \]
Unfortunately, $\pi_2$ is not representable, so there is no forgetting supports morphism $\pi_{2,!} \to \pi_{2,*}$. To circumvent this issue, we will construct an open substack
\begin{equation}
\label{eqn:trans_open}
  j : \left( \Bunc_N \!\underset{\Bun_G}{\times}\! \Bunc_G \right)^\trans \mon \Bunc_N \!\underset{\Bun_G}{\times}\! \Bunc_G
\end{equation}
containing $\Bunc_N \times \BB Z_G$ with the following features:

\begin{proposition}
\label{proposition:cis_vanishing}
Let
\[
   i: \left( \Bunc_N \!\underset{\Bun_G}{\times}\! \Bunc_G \right)^\cis \mon \Bunc_N \!\underset{\Bun_G}{\times}\! \Bunc_G
\]
denote the closed complement to $j$. Then
  \[\pi_{2,*} \cdot i_* \cdot i^! \left( \WW \shoximes r_{G,!} \right) \simeq 0.\]
\end{proposition}

\begin{proposition}
\label{proposition:trans_rep}
The projection $\pi_2 \cdot j$ onto $\Bun_G$ is representable.
\end{proposition}

\subsubsection{} We think of this open substack a locus of transversality, and we will write $(-)^\trans$ for restriction along $j$. Then $\eta_G$ is defined via the diagram:
\[
\begin{tikzcd}[
  column sep={6em,between origins},
  row sep={4em,between origins}
]
  & \pi_{2, !} \cdot \left( \WW \stoximes r_! \right)^\trans \arrow[dr] \arrow[dl, "\sim"] & & \pi_{2, *} \cdot \left( \WW \shoximes r_! \right) [2 \delta_G] \arrow[dl, "\sim"] \\
  \pi_{2, !} \cdot \left( \WW \stoximes r_! \right) & & \pi_{2, *} \cdot \left( \WW \shoximes r_! \right)^\trans [2 \delta_G] &
\end{tikzcd}
\]
The left diagonal map is an isomorphism because $\WW \stoximes r_!$ is $!$-extended from $\Bun_N \times \BB Z_G$. The right diagonal map is an isomorphism by virtue of Proposition \ref{proposition:cis_vanishing}. The middle diagonal map is a combination of the purity morphism and the forgetting supports morphism for $\pi_2 \cdot j$, which makes sense by Proposition \ref{proposition:trans_rep}.

\subsubsection{} We now give the definition of (\ref{eqn:trans_open}). The proofs of Propositions \ref{proposition:cis_vanishing} and \ref{proposition:trans_rep} will appear in the subsequent subsections.

\begin{proposition}
\label{proposition:open_orbit}
The $N \times G$ orbit of $s$ is an open subscheme $\Vin_G^{\circ\circ} \mon \Vin_G^\circ$. It is equal to the $B \times G$ orbit of $s$.
\end{proposition}

\subsubsection{Proof of Proposition \ref{proposition:open_orbit}} Since smooth morphisms are open, it suffices to show that the morphism
  \[N \times T_\ad^+ \times G \to N \times \Vin_G^\circ \times G \to \Vin_G^\circ\]
given by acting on the unit section is smooth. This is a map of schemes smooth over $T_\ad^+$, so we can check the smoothness fiberwise using the explicit description (\ref{eqn:Vin_G_strata}) of the fibers.

\subsubsection{} Now
    \[
    \Bunc_B \underset{\Bun_G}{\times} \Bunc_G
    \simeq \Maps^\gen\left(X, B \bs \Vin_G^\circ / G \mon T \bs \overline{N\bs G} \overset{G}{\times} \Vin_G / G \right)
    \]
and we define
    \[
    \left( \Bunc_B \underset{\Bun_G}{\times} \Bunc_G \right)^\trans
    = \Maps^\gen\left(X, B \bs \Vin_G^{\circ\circ} / G \mon T \bs \overline{N\bs G} \overset{G}{\times} \Vin_G / G \right)
    \]
We take
    \[
    \left( \Bunc_N \underset{\Bun_G}{\times} \Bunc_G \right)^\trans
    = \Bunc_N \underset{\Bunc_B}{\times} \left( \Bunc_B \underset{\Bun_G}{\times} \Bunc_G \right)^\trans
    \]
The map
    \[
    \Bunc_N \times \BB Z_G \to \Bunc_N \underset{\Bun_G}{\times} \Bunc_G
    \]
factors through the transverse locus because $\Vin_G^{\circ\circ}$ contains $\Vin_G^\circ \underset{T^+_\ad}{\times} T_\ad$.

\subsubsection{Sanity check} Let us study the extent to which this device depends on the choice of open substack. Suppose that we have a diagram
\[
\begin{tikzcd}
  \UU \arrow[r, hook, "j"] & \YY \arrow[d, "\pi"] \\
  & S
\end{tikzcd}
\]
in which $j$ is an open embedding and $\pi$ is representable. Then for any map $\FF \to \GG$ of sheaves on $\YY$, we have an obvious commutative diagram:
\[
\begin{tikzcd}[
  column sep={5em,between origins},
  row sep={4em,between origins}
]
  j_! \cdot \FF_\UU \arrow[r] \arrow[d] & \FF \arrow[r] \arrow[d] & j_* \cdot \FF_\UU \arrow[d] \\
  j_! \cdot \GG_\UU \arrow[r] & \GG \arrow[r] & j_* \cdot \GG_\UU
\end{tikzcd}
\]
This induces
\[
\begin{tikzcd}[
  column sep={5em,between origins},
  row sep={1.5em,between origins}
]
  \pi_! \cdot j_! \cdot \FF_\UU \arrow[r] \arrow[dd] & \pi_! \cdot \FF \arrow[dr] \arrow[ddd] & \\
   & & \pi_* \cdot j_* \cdot \FF_\UU \arrow[dd] \\
  \pi_! \cdot j_! \cdot \GG_\UU \arrow[dr] & & \\
   & \pi_* \cdot \GG \arrow[r] & \pi_* \cdot j_* \cdot \GG_\UU
\end{tikzcd}
\]
This diagram exhibits the sense in which our morphism does not depend on the choice of $\UU$. 

\subsubsection{} This sanity check shows that the definition of $\eta_G$ would be unchanged if we replaced $\Bunc_N$ with $\Bun_N$ everywhere.

\subsection{Example: what happens for $\SL_2$}

\subsubsection{} The rest of this paper has no logical dependencies on this subsection.

\subsubsection{} We fix $G = \SL_2$ for this subsection. In this case the objects at hand admit the following descriptions:
\begin{itemize}
    \item The stack $\Vin\Bun_G$ classifies $(\EE_1, \EE_2, \alpha)$ where the $\EE_i$ are a pair of $\SL_2$-bundles and $\alpha$ is a nonzero map of coherent sheaves
        \[\alpha : \EE_1 \to \EE_2.\]
    The torus $T \simeq \Gm$ acts by scaling $\alpha$.
    \item The quotient $\Bunc_G$ classifies $(\EE_1, \EE_2, L, \alpha)$ where the $\EE_i$ are a pair of $\SL_2$-bundles, $L$ is a line, and $\alpha$ is a nonzero map of coherent sheaves
        \[\alpha : \EE_1 \to \EE_2 \otimes L.\]
    The open stratum $\Bunc_{G,G}$ is the locus where $\alpha$ is an isomorphism, and the closed stratum $\Bunc_{G,B}$ is the locus where $\alpha$ has rank one.
    \item The fiber product $\Bunc_N \underset{\Bun_G}{\times} \Bunc_G$ classifies the data of
        \[\omega^{1/2} \mon \EE_1 \to \EE_2 \otimes L.\]
    The transverse locus is the open substack where the composed map $\omega^{1/2} \to \EE_2 \otimes L$ is nonzero.
\end{itemize}

\subsubsection{} Let's verify by hand that the transverse locus is representable over $\Bun_G$. The $G$-stratum of the transverse locus is simply
    \[\Bunc_N \times \BB Z_G \to \Bun_G\]
so we may restrict attention to the $B$-stratum. This stratum has a further stratification by the degree of $\im(\alpha)$. Now the stratum where $\im(\alpha)$ has degree $d$ can be perceived as classifying a point $\EE_2$ of $\Bun_G$ along with:
\begin{itemize}
    \item[(i)] a choice of degree $d$ divisor $D$,
    \item[(ii)] a line $L$ along with an injective map of coherent sheaves
        \[\omega^{1/2}(D) \mon \EE_2 \otimes L,\]
    and
    \item[(iii)] an $\SL_2$ bundle $\EE_1$ fitting into a triangle 
        \[
        \begin{tikzcd}[
          column sep={1em},
          row sep={1em}
        ]
            \omega^{1/2} \arrow[d, hook] \arrow[dr, hook] & \\
            \EE_1 \arrow[r, ->>] & \omega^{1/2}(D)
        \end{tikzcd}
        \]
\end{itemize}
We claim that the stacks parameterizing each additional data are schematic over the previous ones. So for (i) we are looking at the projection
    \[\Sym^d(X) \times \Bun_G \to \Bun_G,\]
which is schematic because $\Sym^d(X)$ is a scheme. Then the choices in (ii) are parameterized by the generalized projective bundle over $\Sym^d(X) \times \Bun_G$ whose fibers are 
    \[\bfP H^0\left( \Maps(\omega^{1/2}(D), \EE_1) \right).\]
Finally, the data in (iii) are parametrized by the vector bundle over this projective bundle whose fibers are
\begin{align*}
    & \fib \left\{ \Maps\left(\omega^{1/2}(D), \omega^{-1/2}(-D)[1] \right) \to \Maps\left( \omega^{1/2}, \omega^{-1/2}(-D)[1] \right) \right\}   \\
    & \qquad \qquad \simeq \Maps\left(\omega^{1/2}_D, \omega^{-1/2}(-D) [1] \right)
    \simeq \Maps\left( \OO, (\omega_D)^{\otimes 2} \right)^\vee.
\end{align*}

\subsection{Proof of Proposition \ref{proposition:cis_vanishing}}

\subsubsection{} We prove Proposition \ref{proposition:cis_vanishing} by a series of reductions. Let us write
    \[\FF = \WW \shoximes r_{G,!}.\]

\subsubsection{First reduction} We use the stratification of $\Bunc_G$ into the strata $\Bunc_{G,P}$. It suffices to show that the !-restrictions of $\FF$ to each
  \[\left(\Bunc_N \underset{\Bun_G}{\times} \Bunc_{G,P}\right)^\cis\]
are killed by $*$-pushforward to $\Bun_G$.

\subsubsection{A generality} Let $T$ be any algebraic torus, and consider a principal $T$-bundle
\[
\begin{tikzcd}[
  column sep={4em,between origins},
  row sep={3em,between origins}
]
    P \arrow[d, "f" left] \arrow[r] & \pt \arrow[d] \\
    S \arrow[r, "\alpha"] & \BB T
\end{tikzcd}
\]
For any sheaf $\FF$ on $S$ we have
    \[\Gamma(S, \FF) \simeq \Maps(k_{\BB T}, \alpha_* \cdot \FF) \simeq \Maps_{\rH_*(T)}\left(r^! \cdot k_{\BB T}, \Gamma(P, \FF) \right).\]

\subsubsection{Second reduction} Let $\Vin\Bun_{G,P}$ denote the fiber of $\Vin\Bun_G$ over the point $c_P \in T_\ad^+$. We have
    \[\Bun_{G,P} \simeq \Vin\Bun_{G,P}/Z_M.\]
By the above, it suffices to show that the !-pullbacks of $\FF$ to
  \[\left( \Bunc_N \underset{\Bun_G}{\times} \VBun_{G,P} \right)^\cis\]
are killed by $*$-pushforward to $\Bun_G$.

\subsubsection{Third reduction} Recall that $\VBun_{G,P}$ has a stratification of the form
\begin{equation}
\label{eqn:VinBun_G_str}
    \Bun_{P^-} \underset{\Bun_M}{\times} \HH_M^+ \underset{\Bun_M}{\times} \Bun_P \to \VBun_{G,P}.      
\end{equation}
Therefore, it suffices to show that the !-pullback of $\FF$ to this stratification is killed by $*$-pushforward along
  \[\left( \Bunc_N \underset{\Bun_G}{\times} \Bun_{P^-} \underset{\Bun_M}{\times} \HH_M^+ \underset{\Bun_M}{\times} \Bun_P \right)^\cis \to \Bun_P \to \Bun_G.\]

\subsubsection{Fourth reduction} Finally, we recall that  
  \[\Bunc_N \underset{\Bun_G}{\times} \Bun_{P^-}\]
has a relative position stratification by
  \[W/W_P.\]
It suffices to prove the analogous statement for
  \[
  \pi_2^\gr : \left( \Bunc_N \underset{\Bun_G}{\times} \Bun_{P^-} \right)^w \underset{\Bun_M}{\times} \HH_M^+ \underset{\Bun_M}{\times} \Bun_P  \to \HH_M^+ \underset{\Bun_M}{\times} \Bun_P
  \]
where $w \neq 1$.

\subsubsection{End of the proof} By \cite[Lemma 4.2.1]{lin_nearby_cycles}, the !-restriction of $r_!$ to the above stack is !-pulled back along $\pi_2^\gr$. Therefore, it suffices to observe that $\WW$ is annihilated by $*$-pushforward along
    \[
    \left( \Bunc_N \underset{\Bun_G}{\times} \Bun_{P^-} \right)^w 
    \to \Bun_M.
    \]
This is a standard fact whose proof we review in \S\ref{subsection:Jacquet_Whittaker}.

\subsection{Proof of Proposition \ref{proposition:trans_rep}}

\subsubsection{} To simplify notation, we will prove the analogue of Proposition \ref{proposition:trans_rep} for the \emph{untwisted} version of $\Bun_N$. The twisted version can be proved in the same way by using the appropriate non-constant group scheme twist of $N$ and replacing the notion of generic maps with generic sections.

\subsubsection{Proof of Proposition \ref{proposition:trans_rep}} Recall that an Artin stack is an algebraic space if the stabilizer group of every geometric point is trivial. Therefore, we only need to check that our map induces injections on stabilizer groups when evaluated at geometric points (note that this property is stable under base change). We shall do this after passage to a stratification

\subsubsection{Stratification} We will again use the stratification of $\Bunc_G$ into $\Bunc_{G,P}$. We further stratify $\Bunc_{G,P}$ using the quotient of (\ref{eqn:VinBun_G_str}) by $Z_M$. Therefore, it suffices to show that the map
    \[
    \left(\Bunc_N \underset{\Bun_G}{\times} \Bun_{P^-} \right)^\trans \underset{\Bun_M}{\times} \HH_M^+/Z_M \underset{\Bun_M}{\times} \Bun_P \to \Bun_G
    \]
is representable. Since the morphism $\Bun_P \to \Bun_G$ is schematic, we only need to check that the map
    \[
    \alpha : \left(\Bunc_N \underset{\Bun_G}{\times} \Bun_{P^-} \right)^\trans \underset{\Bun_M}{\times} \HH_M^+/Z_M \to \Bun_M
    \]
is representable. 

\subsubsection{} Now recall that the source of the stratification (\ref{eqn:VinBun_G_str}) is given by
    \[
    \Maps^\gen\left(X, P^- \bs M/P \mon P^- \bs M^+ /P \right)
    \]
for a certain monoid $M^+$ containing $M$ as its unit group. Therefore, in terms of generic maps, the morphism $\alpha$ can be written as the composite:
\begin{align*}
    & \Maps^\gen\left( X, \BB N_M \mon N \bs G \overset{P^-}{\times} M^+/M \right)/Z_M \\
    & \qquad \mon \Maps^\gen\left( X, \BB (N_M \cdot Z_M) \mon N \bs G \overset{P^-}{\times} M^+/M \cdot Z_M \right) \\
    & \qquad \qquad \overset{\beta}{\longrightarrow} \Maps(X, \BB M)
\end{align*}
The first arrow is a closed immersion because it is a base change of $\BB M \to \Bun_M$. Therefore, it suffices that to show that $\beta$ induces injections on stabilizer groups when evaluated at geometric points. 

\subsubsection{} Let us write
    \[
    \YY = N \bs G \overset{P^-}{\times} M^+/Z_M \cdot M
    \qquad
    V = N_M \cdot Z_M
    \]
for convenience. Then $\beta$ fits into the following Cartesian diagram. 
\[
\begin{tikzcd}[
  column sep={12em,between origins},
  row sep={4em,between origins}
]
    \Maps^\gen\left(X, \BB V \mon \YY \right) \arrow[r] \arrow[d] & \Maps(X, \YY) \arrow[d, "\star"] & \\
    \Maps^\gen\left(X, \BB V \to \BB M \right) \arrow[r] \arrow[d] & \Maps^\gen\left(X, \YY \to \BB M \right) \arrow[r] \arrow[d] & \Maps(X, \BB M) \arrow[d] \\
    \Maps^\gen\left(X, \BB V \right) \arrow[r] & \Maps^\gen\left(X, \YY \right) \arrow[r] & \Maps^\gen(X, \BB M)
\end{tikzcd}
\]
The composite of the lower row induces injections on stabilizers of geometric points, and hence so does the composite of the middle row. Therefore, it suffices to show that the map ($\star$) has the same property. But in fact this property holds for the composite
    \[\Maps(X, \YY) \to \Maps^\gen(X, \YY)\]
because $\YY$ is the quotient of a scheme by an algebraic group.